\section{A second decision rule: conformal open-set detection}\label{sec:conf}
The real-data replay left one hypothesis unmet: we could not show that the open-world planner is no less accurate on ordinary supernovae (H5). A second, independently preregistered study attacked that cost directly, with two changes on top of the causal configuration of Section~\ref{sec:real}.

\paragraph{Method.} \emph{Conformal decision.} After any planner's measurements $D$, compute the prior-free score $s=\log p(D\mid h_0)-\max_c\log p(D\mid c)$ and flag ``unknown'' when $s>\tau$, where $\tau$ is the largest, over known classes, of the split-conformal $(1-\alpha)$ quantile of $s$ on calibration events of that class. If calibration and new events of a class are exchangeable, the probability of a false ``unknown'' is at most $\alpha$ for every known class, for any planner and any class mix \citep{Vovk2005,Angelopoulos2023}; otherwise the closed-world decision stands. \emph{Lookahead.} Because greedy one-step utilities end with too few measurements, a planner L adds $\lambda$ times the value of the best remaining later epochs to P's utility ($\lambda=0$ recovers P).

\paragraph{Protocol.} $\lambda=\ConfLambda$ was chosen on the validation split; $\alpha=\ConfAlpha$ was fixed. Calibration and sealed known-class test sets were drawn fresh from Bright Transient Survey supernovae never used by any earlier analysis, split by a hash of the object name: \ConfNCal\ calibration and \ConfNFresh\ sealed test supernovae, the latter loaded for the first time in the test run. No unused unknown-class events remain, so the unknown-detection hypotheses (HC4--HC6) reuse the test split of Section~\ref{sec:real}, an additional declared look. Each of six hypotheses had to pass on all three noise seeds.

\begin{table*}[t]
\centering\scriptsize
\caption{Conformal open-set study, preregistered results per noise seed: estimate [95\% CI]; $^\times$ = criterion not met. Primary system: lookahead planner L with the conformal decision (C-L). $\Delta$: C-L minus closed-world EIG (B3), except HC6 (AUROC of $s$, L minus P). HC1 uses Clopper--Pearson intervals; HC2--HC4 and HC6 paired bootstrap.}
\label{tab:conf}
\begin{tabular}{llll}
\toprule
Hypothesis & Seed 1 & Seed 2 & Seed 3 \\
\midrule
HC1 false ``unknown'', sealed SNe ($\leq\alpha$) & +0.035 [+0.020, +0.055] & +0.061 [+0.042, +0.086] & +0.063 [+0.043, +0.089] \\
HC2 $\Delta$acc.\ SNe ($>-0.05$) & +0.082 [+0.039, +0.123] & +0.074 [+0.033, +0.115] & +0.059 [+0.012, +0.106] \\
HC3 $\Delta$CW SNe ($<+0.05$) & -0.027 [-0.057, +0.002] & -0.049 [-0.078, -0.020] & -0.012 [-0.043, +0.018] \\
HC4 $\Delta$CW unknowns ($<0$) & -0.239 [-0.375, -0.102] & -0.261 [-0.409, -0.114] & -0.193 [-0.341, -0.045] \\
HC5 unknowns flagged ($>\alpha$) & +0.420 [+0.316, +0.530] & +0.545 [+0.436, +0.652] & +0.432 [+0.327, +0.542] \\
HC6 $\Delta$AUROC, L$-$P ($>0$) & +0.012 [-0.013, +0.038]$^\times$ & +0.031 [+0.006, +0.058] & +0.039 [+0.005, +0.075] \\
\bottomrule
\end{tabular}

\end{table*}

\paragraph{Results.} Five of six hypotheses passed (Table~\ref{tab:conf}). On the sealed supernovae, the false-``unknown'' rate was \ConfFa, within the certified \ConfAlphaPct\ (HC1 \ConfHCOneVerdict), and C-L was \ConfAccGain\ \emph{more} accurate than closed-world EIG (HC2 \ConfHCTwoVerdict, the re-test of H5) with no more confident errors (HC3 \ConfHCThreeVerdict). On the real unknowns, C-L flagged \ConfFlag\ (HC5 \ConfHCFiveVerdict) and was less often confidently wrong than B3 (HC4 \ConfHCFourVerdict). Lookahead did not reliably improve the ranking of unknowns over P (HC6 \ConfHCSixVerdict; AUROC difference \ConfAucGain, not significant on seed 1).

\paragraph{Where the gains come from, and what they cost.} With causal measurements, P without the conformal layer was already as accurate as B3 on the sealed supernovae (\ConfAccFreshP\ against \ConfAccFreshBthree), with a raw false-``unknown'' rate of \ConfRawFaFreshP; much of the earlier H5 failure was the retrospective protocol. Lookahead added accuracy (\ConfAccFreshCL) by using more of the budget (\ConfObsL\ measurements against \ConfObsP\ for P). The conformal layer is planner-agnostic: applied to closed-world EIG, it flagged \ConfFlagBthreeConf\ of unknowns while keeping B3's supernova accuracy near its own (\ConfAccFreshBthreeConf). The cost is on the unknowns. C-L was confidently wrong on \ConfCwUnkCL\ of them, against \ConfCwUnkP\ for P's Bayesian decision and \ConfCwUnkBthree\ for B3. The conformal rule buys a certified false-alarm rate and accuracy on known classes; the Bayesian $h_0$ buys fewer confident errors on unknowns, with no guarantee. Which to deploy depends on which error is more expensive.
